\documentclass[11pt, a4paper]{article}
\usepackage[margin=1in]{geometry}

\title{Code Typesetting in rushtex}
\author{rushtex team}
\date{September 2026}

\begin{document}
\maketitle

\section{Overview}
The \texttt{rushtex} typesetting engine now natively supports raw code listings, verbatim environments, and inline code snippets using the standard PDF \texttt{Courier} monospace typeface family.

\section{Raw Verbatim Block}
Standard raw verbatim syntax with unescaped symbols like \verb|{}|, \verb|%|, and \verb|\|:

\begin{verbatim}
fn main() {
    // Everything inside verbatim is literal
    let message = "Hello from rushtex!";
    println!("{}", message);
}
\end{verbatim}

\section{Numbered Code Listing}
Code listings with line numbers using the \texttt{lstlisting} syntax:

\begin{lstlisting}[language=Rust, numbers=left]
pub struct Vector3 {
    pub x: f32,
    pub y: f32,
    pub z: f32,
}

impl Vector3 {
    pub fn dot(&self, other: &Self) -> f32 {
        self.x * other.x + self.y * other.y + self.z * other.z
    }
}
\end{lstlisting}

\section{Inline Code Snippets}
You can use \texttt{cargo build --release} or \verb|rushtex input.tex -o output.pdf| directly in your terminal.

\end{document}
